\begin{abstract}
Traditional Multi-Factor Authentication (MFA) mechanisms fail when the client endpoint is untrusted or compromised by infostealer malware, screen grabbers, or adversary-in-the-middle proxies. While hardware-bound tokens (FIDO2/WebAuthn) protect against network relay, they remain vulnerable to malware-induced session hijacking post-authentication and face physical availability constraints in ad-hoc web and mobile environments. This paper proposes a Blind Cognitive Multi-Factor Authentication (BC-MFA) architecture designed for untrusted web and mobile public clients. The protocol decouples the visual display on the terminal from the user's authentic credential computation through dynamic, quadrant-balanced toroidal challenge-response matrices, paired with passive behavioral anomaly detection (cognitive hesitation, pointer/touch kinematics, and keystroke cadence). We formalize the protocol exchange under OpenID Connect (OIDC) Core 1.0 and OAuth 2.0, combining RFC 7636 PKCE, registered RFC 8176 Authentication Method Reference (\texttt{amr}) claims, and RFC 9449 Demonstrating Proof-of-Possession (DPoP) access token binding. Analytical threat modeling and a reproducible 500-trial Monte Carlo intersection simulation quantify resilience against keyloggers, screen scrapers, SIM swapping, multi-session intersection attacks, and automated relay bots without requiring dedicated client cryptographic hardware, while explicitly scoping first-contact active reverse-proxy limits. Finally, we establish a literature-grounded latency model and the formal empirical evaluation plan for the Master's dissertation.
\end{abstract}

\begin{IEEEkeywords}
\normalfont
Cognitive authentication, untrusted endpoints, infostealers, OpenID Connect, RFC 8176, RFC 9449, behavioral anomaly detection, challenge-response.
\end{IEEEkeywords}
